\documentclass[12pt]{article}
\usepackage[utf8]{inputenc}
\usepackage[T1]{fontenc}
\usepackage[brazil]{babel}
\usepackage[a4paper,margin=2.5cm]{geometry}
\usepackage{enumitem}
\setlist[itemize]{noitemsep,topsep=4pt}
\title{Sprint 3 - S3\_02\\Paradigmas: Recursão e Funções de Alta Ordem}
\author{Matheus Moreli Durão Silva\\Matrícula 12621CBS217\\Cibersegurança - Bacharelado - Noturno}
\date{2026}
\begin{document}
\maketitle

\section*{Entregável}
Algoritmos recursivos e comparação com abordagens iterativas e funcionais. Este documento registra o objetivo do card, a solução desenvolvida e a forma de conferência. Os arquivos de código e as evidências citadas permanecem na mesma pasta deste documento.

\section*{O que foi estudado e aplicado}
\begin{itemize}
\item Caso-base e chamada recursiva
\item Fibonacci com memoização
\item Uso de lambda, map e filter
\item Comparação de resultados e limites da recursão
\end{itemize}

\section*{Como conferir}
O avaliador pode abrir os arquivos desta pasta, executar os exemplos quando houver código e comparar a saída com a explicação apresentada. Nos cards ligados ao Inventário de Ativos, a solução foi relacionada ao Trabalho 2 sem modificar requisitos já atendidos do CRUD.

\section*{Comparação medida}
O script \texttt{comparar\_desempenho.py} registrou a mediana de cinco execuções em Python 3.13.15/Windows. Para fatorial de 800, iteração levou 0,074453 ms/chamada e recursão 0,114057 ms/chamada; ambas têm tempo $O(n)$, mas a recursão acrescenta pilha $O(n)$. Para Fibonacci de 30, a versão sem cache levou 66,0056 ms/chamada e a memoizada 0,0037 ms/chamada. Sem cache há recomputação de subproblemas; com cache, tempo e memória são $O(n)$.
\section*{Funções de alta ordem}
\texttt{map} extrai nomes, \texttt{filter} seleciona ativos de criticidade Alta/Crítica, \texttt{reduce} soma comprimentos e \texttt{sorted} com \texttt{lambda} ordena por nome. A tabela completa, o método de medição e suas limitações constam em \texttt{analise\_comparativa.md}; os valores absolutos podem mudar em outra máquina.
\section*{Conclusão}
Os exemplos mostram que recursão pode ser clara, mas memoização é decisiva em problemas com subproblemas repetidos.
\end{document}
